\chapter{Sistem Real-Time dan Asinkron}
\label{ch:realtime}

\section*{Tujuan Pembelajaran}
\begin{itemize}
  \item Mengukur keterlambatan dan biaya \textit{polling}, SSE (Server-Sent
        Events), serta WebSocket pada peristiwa yang sama, lalu memilih satu
        di antaranya beserta alasannya.
  \item Merancang alur asinkron memakai antrean, \textit{worker}, dan
        \textit{webhook}, lalu menunjukkan jeda \textit{retry}-nya pada hasil
        pengukuran.
  \item Membedakan pengiriman \textit{at-least-once} dari pengiriman sekali
        saja, lalu membuktikan \textit{consumer} yang idempoten menahan
        duplikatnya.
\end{itemize}

\section{Pendahuluan}

Dua kebutuhan sering disatukan padahal berbeda. Kebutuhan pertama, pengguna
ingin melihat perubahan segera, tanpa menekan tombol \textit{refresh}.
Kebutuhan kedua, pekerjaan yang lama tidak boleh menahan jawaban, sehingga
harus dikerjakan di belakang.

Keduanya menolak pola permintaan dan jawaban yang biasa. Yang pertama
menuntut jalan dari server ke klien, padahal HTTP dirancang untuk arah
sebaliknya. Yang kedua menuntut pekerjaan hidup lebih lama daripada
permintaan yang memulainya, sehingga harus dititipkan ke tempat lain.

Demonya dua bagian yang berjalan bersama: satu server peristiwa yang
melayani tiga cara pengiriman, dan satu alur asinkron berisi antrean,
\textit{worker}, dan penerima \textit{webhook}. Referensi utamanya adalah
HTML Standard untuk SSE beserta \texttt{EventSource} \cite{whatwg-html}, RFC
6455 untuk WebSocket \cite{rfc6455}, RFC 9110 untuk arti metode yang
idempoten \cite{rfc9110}, dan dokumentasi Amazon SQS untuk
\textit{visibility timeout} beserta akibatnya \cite{aws-sqs-visibility}.

\section{Polling, Server-Sent Events, dan WebSocket}

Tiga cara dipakai luas untuk membawa peristiwa dari server ke klien.
\textit{Polling} membuat klien bertanya berulang kali, dan jawabannya sering
kosong. SSE membuat satu sambungan HTTP tetap terbuka, lalu server menulis
peristiwa ke sambungan itu \cite{whatwg-html}. WebSocket menukar sambungan
HTTP menjadi sambungan dua arah lewat jabat tangan khusus \cite{rfc6455}.

Analoginya seperti menunggu nilai ujian keluar. \textit{Polling} berarti
menengok papan pengumuman tiap lima menit. SSE berarti menunggu di depan
papan sementara petugas menempel pengumuman begitu siap. WebSocket berarti
berdiri di depan petugas, sehingga keduanya dapat saling berbicara.

Gambar \ref{fig:transport} memperlihatkan ketiganya pada demo bab ini.

\begin{figure}[!htb]
  \centering
  \begin{tikzpicture}[
    font=\scriptsize,
    kotak/.style={draw, rounded corners, minimum height=0.8cm,
                  minimum width=2.2cm, align=center, font=\scriptsize\bfseries},
    panah/.style={-{Stealth[length=2mm]}, thick}
  ]
    \node[kotak, fill=praditagreen!15, minimum height=2.6cm] (server) at (0,0) {Server\\peristiwa\\{\tiny\mdseries port 8100}};
    \node[kotak, fill=gray!35] (polling) at (4.3,1.7) {Polling\\{\tiny\mdseries 11 permintaan}};
    \node[kotak, fill=blue!8] (sse) at (4.3,0) {SSE\\{\tiny\mdseries 1 sambungan}};
    \node[kotak, fill=orange!12] (ws) at (4.3,-1.7) {WebSocket\\{\tiny\mdseries 1 sambungan}};

    \draw[panah] (polling.west) -- node[font=\tiny, above] {tanya} (server.east |- polling);
    \draw[panah] (server.east |- polling) ++(0,-0.3) -- ++(2.1,0);
    \draw[panah] (server.east |- sse) -- node[font=\tiny, above] {dorong} (sse.west);
    \draw[panah] (server.east |- ws) -- node[font=\tiny, above] {dorong} (ws.west);
    \draw[panah] (ws.west) ++(0,-0.3) -- ++(-2.1,0);
    \node[font=\tiny, gray, anchor=west] at (6.0,1.4) {jawaban sering kosong};
    \node[font=\tiny, gray, anchor=west] at (6.0,0) {satu arah};
    \node[font=\tiny, gray, anchor=west] at (6.0,-2.0) {dua arah};
  \end{tikzpicture}
  \caption{Tiga cara membawa peristiwa ke klien, beserta jumlah permintaan
           yang terukur untuk 10 peristiwa}
  \label{fig:transport}
\end{figure}

Listing \ref{lst:sse13} memperlihatkan inti SSE di sisi server. Panjang
isinya tidak disebutkan, karena sambungannya belum berakhir, dan tiap
peristiwa ditulis sebagai satu blok yang diakhiri baris kosong.

\begin{lstlisting}[style=PythonStyle, caption={Satu sambungan yang tetap terbuka, tersedia di \berkas{sources/chapter13/aplikasi.py}}, label={lst:sse13}]
    self.send_response(200)
    self.send_header("Content-Type", JENIS_ISI_SSE)
    self.send_header("Cache-Control", "no-store")
    self.send_header("Connection", "keep-alive")
    self.end_headers()
    try:
      while True:
        try:
          peristiwa = antrean.get(timeout=JEDA_PENANDA_HIDUP_DETIK)
        except queue.Empty:
          self.wfile.write(b": tetap hidup\n\n")
          self.wfile.flush()
          continue
        blok = (f"id: {peristiwa['nomor']}\n"
                f"data: {json.dumps(peristiwa['isi'])}\n\n")
        self.wfile.write(blok.encode())
        self.wfile.flush()
\end{lstlisting}

Baris penanda hidup itu bukan hiasan. Sambungan yang lama menganggur dapat
ditutup perantara di tengah jalan, dan kliennya baru tahu ketika peristiwa
berikutnya tidak pernah tiba. Satu baris komentar tiap 15 detik cukup untuk
mencegahnya.

Tabel \ref{tab:transport} memuat hasil pengukuran ketiganya. Sepuluh
peristiwa yang sama diterbitkan untuk tiap transport, lalu keterlambatan,
jumlah byte yang diterima klien, dan jumlah permintaan HTTP dicatat.

\begin{table}[!htb]
  \centering
  \caption{Tiga transport pada 10 peristiwa yang sama, hasil pengukuran
           nyata}
  \label{tab:transport}
  \begin{tabularx}{\textwidth}{Xllll}
    \hline
    \textbf{Transport} & \textbf{Tengah ms} & \textbf{Tertinggi ms} & \textbf{Byte} & \textbf{Permintaan} \\
    \hline
    \textit{polling}, jeda 0,5 detik & 199,8 & 202,0 & 639 & 11 \\
    \hline
    SSE & 1,8 & 2,0 & 400 & 1 \\
    \hline
    WebSocket & 1,9 & 2,1 & 490 & 1 \\
    \hline
  \end{tabularx}
\end{table}

Tabel itu punya dua konsekuensi praktis. Pertama, keterlambatan
\textit{polling} ditentukan jedanya, bukan kecepatan servernya: 199,8
milidetik pada jeda 0,5 detik, sedangkan kedua transport yang mendorong
berada di bawah 2 milidetik. Kedua, satu sambungan mengganti sebelas
permintaan, dan sebelas permintaan itu masing-masing membawa header
permintaan dan header jawaban yang tidak ikut terhitung di kolom byte.

Kolom byte perlu dibaca hati-hati. Yang dihitung hanya isi yang diterima
klien, yaitu badan jawaban untuk \textit{polling}, baris data untuk SSE, dan
isi frame untuk WebSocket. SSE mencapai 400 byte karena tiap peristiwa
ditulis sebagai dua baris teks, sedangkan WebSocket mencapai 490 byte karena
tiap frame membawa seluruh peristiwa sebagai JSON (JavaScript Object
Notation) ditambah dua byte header. Selisih 90 byte itu berasal dari isi yang
dibawa, bukan dari protokolnya.

Perbedaan yang menentukan pilihan justru tidak ada di tabel. SSE berjalan
satu arah dan otomatis menyambung ulang di browser, sedangkan WebSocket
berjalan dua arah tetapi penyambungan ulangnya ditulis sendiri. Karena itu
papan pemberitahuan cukup dengan SSE, sedangkan permainan atau penyuntingan
bersama membutuhkan WebSocket.

\section{Publish/Subscribe dan Satu Sumber Peristiwa}

\textit{Publish/subscribe} memisahkan yang menerbitkan dari yang membaca.
Penerbit tidak tahu ada berapa pembaca, dan pembaca tidak tahu siapa yang
menerbitkan. Pada demo bab ini, pemisahan itu membuat ketiga transport
membaca peristiwa yang sama, sehingga perbedaan yang terukur memang berasal
dari cara pengirimannya.

Listing \ref{lst:terbit13} memperlihatkan penerbitannya. Dua hal dikerjakan
di dalam satu kunci: pemberian nomor dan penitipan ke tiap pelanggan.

\begin{lstlisting}[style=PythonStyle, caption={Satu penerbitan, banyak pelanggan, tersedia di \berkas{sources/chapter13/aplikasi.py}}, label={lst:terbit13}]
def terbitkan(isi):
  """Mencatat satu peristiwa, lalu menitipkannya ke tiap pelanggan SSE.

  Pencatatan dan penitipan dilakukan di dalam satu kunci, agar nomor
  peristiwa tidak pernah terlewat atau terbalik ketika dua penerbit datang
  bersamaan.
  """
  with KUNCI:
    peristiwa = {"nomor": len(RIWAYAT) + 1, "isi": isi}
    RIWAYAT.append(peristiwa)
    for antrean in list(PELANGGAN):
      if antrean.qsize() < BATAS_ANTREAN_PELANGGAN:
        antrean.put(peristiwa)
  return peristiwa
\end{lstlisting}

Batas antrean pada baris kedua dari bawah menjawab pertanyaan yang muncul
berikutnya: bagaimana bila ada pelanggan yang membaca lebih lambat daripada
kecepatan penerbitan. Tanpa batas, server menyimpan peristiwa untuk pelanggan
itu sampai memorinya habis. Dengan batas, pelanggan yang tidak sanggup
mengikuti kehilangan peristiwa, dan kehilangan itu dipilih dengan sadar.

Nomor peristiwa punya guna kedua. Klien \textit{polling} memakainya sebagai
penanda sampai mana ia sudah membaca, dan klien SSE menerimanya sebagai field
\texttt{id}, yang dipakai browser untuk melanjutkan dari tempat terakhir
ketika sambungannya tersambung ulang.

\section{Antrean, Background Worker, dan Webhook}

Pekerjaan yang lama tidak dikerjakan di dalam permintaan. Permintaan hanya
menitipkannya ke antrean, lalu menjawab bahwa titipannya diterima.
\textit{Worker} yang berjalan terpisah mengambil pekerjaan itu dan
mengerjakannya, misalnya mengirim \textit{webhook} ke sistem lain.

Analoginya seperti loket pengurusan surat. Mahasiswa menyerahkan berkas dan
menerima tanda terima, lalu pulang. Petugas di belakang mengerjakan berkasnya
menurut urutan masuk, dan hasilnya diambil kemudian. Loketnya tidak pernah
menahan orang sampai suratnya jadi.

Gambar \ref{fig:antrean} memperlihatkan alurnya beserta ketiga jenis
pekerjaan yang diukur bab ini.

\begin{figure}[!htb]
  \centering
  \begin{tikzpicture}[
    font=\scriptsize,
    kotak/.style={draw, rounded corners, minimum height=0.85cm,
                  minimum width=2.2cm, align=center, font=\scriptsize\bfseries},
    panah/.style={-{Stealth[length=2mm]}, thick}
  ]
    \node[kotak, fill=blue!8] (kirim) at (0,0) {Pengirim\\{\tiny\mdseries 3 pekerjaan}};
    \node[kotak, fill=praditagreen!15, minimum height=1.1cm] (antrean) at (3.2,0) {Antrean\\{\tiny\mdseries siap\_pada}\\{\tiny\mdseries percobaan}};
    \node[kotak, fill=orange!12, minimum height=1.1cm] (worker) at (6.4,0) {Worker\\{\tiny\mdseries periksa duplikat}\\{\tiny\mdseries lalu kirim}};
    \node[kotak, fill=gray!35] (webhook) at (9.6,0.6) {Webhook\\{\tiny\mdseries port 8102}};
    \node[kotak, fill=red!22] (gagal) at (9.6,-0.9) {500\\{\tiny\mdseries dua kali}};

    \draw[panah] (kirim) -- (antrean);
    \draw[panah] (antrean) -- (worker);
    \draw[panah] (worker.east) -- ++(0.3,0) |- (webhook.west);
    \draw[panah] (worker.east) -- ++(0.3,0) |- (gagal.west);
    \draw[panah] (gagal.south) -- ++(0,-0.45) -| node[font=\tiny, below, pos=0.25] {dicoba lagi dengan jeda berlipat} (antrean.south);
  \end{tikzpicture}
  \caption{Alur asinkron: titipan masuk antrean, \textit{worker}
           mengerjakannya, kegagalan kembali ke antrean}
  \label{fig:antrean}
\end{figure}

Antrean pada demo ini satu file JSON, dan itu memang bukan antrean untuk
produksi. Yang ditirukan hanya sifat yang menentukan cara menulis
\textit{consumer}, dan sifat terpentingnya ada pada Listing
\ref{lst:ambil13}: pekerjaan yang diambil tidak langsung hilang.

\begin{lstlisting}[style=PythonStyle, caption={Pekerjaan yang diambil hanya disembunyikan sementara, tersedia di \berkas{sources/chapter13/antrean.py}}, label={lst:ambil13}]
def ambil(sekarang):
  """Mengambil satu pekerjaan yang sudah waktunya dikerjakan.

  Pekerjaan yang diambil tidak langsung hilang dari antrean. Waktu siapnya
  digeser sejauh jeda backoff, sehingga pekerjaan itu kembali muncul bila
  worker-nya mati sebelum memanggil selesai. Di sanalah sifat at-least-once
  berasal: pekerjaan dapat terkirim lebih dari sekali, dan tidak pernah
  hilang tanpa dikerjakan.
  """
  keadaan = muat()
  for pesan in keadaan["pesan"]:
    if pesan["siap_pada"] <= sekarang:
      pesan["percobaan"] += 1
      pesan["siap_pada"] = sekarang + jeda_backoff(pesan["percobaan"])
      simpan(keadaan)
      return dict(pesan)
  return None
\end{lstlisting}

Pergeseran waktu siap itu disebut \textit{visibility timeout} pada layanan
antrean sungguhan \cite{aws-sqs-visibility}. Akibatnya dua. Pertama,
pekerjaan tidak hilang ketika \textit{worker}-nya mati, karena pekerjaannya
muncul kembali. Kedua, pekerjaan dapat dikerjakan dua kali, karena kematian
itu bisa terjadi sesudah pekerjaannya tuntas tetapi sebelum di-\textit{ack}.

\textit{Webhook} adalah sisi lain dari hal yang sama. Pada \textit{polling},
sistem kita yang bertanya ke sistem lain. Pada \textit{webhook}, sistem lain
yang memanggil alamat kita ketika ada kabar. Keduanya memindahkan masalah ke
tempat yang sama, yaitu apa yang terjadi ketika pemanggilannya gagal.

\section{Retry dan Exponential Backoff}

\textit{Retry} yang tanpa jeda memperburuk keadaan. Tujuan yang sedang
bermasalah justru dibanjiri percobaan baru, tepat ketika ia paling butuh
kesempatan untuk pulih. Karena itu jeda antarpercobaan dibuat berlipat, dan
caranya satu baris, seperti pada Listing \ref{lst:backoff13}.

\begin{lstlisting}[style=PythonStyle, caption={Jeda yang berlipat dua tiap percobaan, tersedia di \berkas{sources/chapter13/antrean.py}}, label={lst:backoff13}]
def jeda_backoff(percobaan):
  """Menghitung jeda sebelum percobaan berikutnya, berlipat dua tiap kali.

  Jeda yang berlipat menahan worker dari membanjiri tujuan yang sedang
  bermasalah. Percobaan pertama memakai jeda dasar, yang kedua dua kali,
  yang ketiga empat kali, dan seterusnya.
  """
  return JEDA_DASAR_DETIK * (2 ** max(0, percobaan - 1))
\end{lstlisting}

Tabel \ref{tab:retry} memuat seluruh percobaan yang dilakukan
\textit{worker} pada satu daur pengukuran. Jeda dasarnya 0,5 detik, dan
pekerjaan \texttt{gagal-dua-kali} ditolak \textit{webhook} dua kali sebelum
diterima.

\begin{table}[!htb]
  \centering
  \caption{Seluruh percobaan \textit{worker} pada satu daur, hasil
           pengukuran nyata}
  \label{tab:retry}
  \begin{tabularx}{\textwidth}{llXl}
    \hline
    \textbf{Pekerjaan} & \textbf{Ke} & \textbf{Hasil} & \textbf{Detik} \\
    \hline
    \texttt{kerja-0001} & 1 & terkirim dan di-\textit{ack} & 0,01 \\
    \hline
    \texttt{kerja-0002} & 1 & gagal, \textit{webhook} 500 & 0,01 \\
    \hline
    \texttt{kerja-0003} & 1 & terkirim, \textit{worker} mati sebelum \textit{ack} & 0,01 \\
    \hline
    \texttt{kerja-0002} & 2 & gagal, \textit{webhook} 500 & 0,51 \\
    \hline
    \texttt{kerja-0003} & 2 & duplikat ditahan & 0,51 \\
    \hline
    \texttt{kerja-0002} & 3 & terkirim dan di-\textit{ack} & 1,52 \\
    \hline
  \end{tabularx}
\end{table}

Kolom detik memperlihatkan jedanya bekerja. Percobaan kedua
\texttt{kerja-0002} terjadi 0,5 detik sesudah yang pertama, dan percobaan
ketiga terjadi 1,0 detik sesudah yang kedua. Angka 0,01 dan 0,51 dan 1,52
adalah waktu sejak pengukuran mulai, sehingga selisihnya persis jeda dasar
lalu dua kali jeda dasar.

Pertanyaan berikutnya, sampai kapan dicoba ulang. Pada demo ini batasnya
lima percobaan, dan pekerjaan yang tetap gagal sesudah itu pantas dipindahkan
ke tempat terpisah, yang biasa disebut \textit{dead letter queue}, agar
antreannya tidak tersumbat satu pekerjaan yang tidak akan pernah berhasil.

\section{Consumer Idempoten dan Pengiriman At-Least-Once}

\textit{At-least-once} berarti satu pekerjaan dijamin terkirim setidaknya
sekali, dan boleh lebih. Jaminan itu bukan kemalasan perancangnya, melainkan
pilihan yang masuk akal: menjamin tepat sekali menuntut kesepakatan antara
antrean dan \textit{consumer} yang mahal, sedangkan menjamin paling banyak
sekali berarti pekerjaan boleh hilang.

Akibatnya pindah ke sisi \textit{consumer}, yang harus menahan duplikat
sendiri. Listing \ref{lst:idempoten13} memperlihatkan urutannya, dan urutan
itulah intinya.

\begin{lstlisting}[style=PythonStyle, caption={Pemeriksaan duplikat dilakukan paling dahulu, tersedia di \berkas{sources/chapter13/worker.py}}, label={lst:idempoten13}]
def kerjakan(pesan, tirukan_crash):
  """Mengerjakan satu pesan, lalu mengembalikan keterangan hasilnya.

  Urutannya penting. Pemeriksaan duplikat dilakukan paling dahulu, sebelum
  webhooknya dikirim, karena itulah satu-satunya tempat duplikat masih dapat
  ditahan tanpa akibat.
  """
  if antrean.sudah_diproses(pesan["id"]):
    antrean.selesai(pesan["id"])
    return "duplikat ditahan"
  status = kirim_webhook(pesan)
  if status != 200:
    return f"gagal, webhook {status}, dicoba lagi nanti"
  antrean.tandai_diproses(pesan["id"])
  if pesan["jenis"] == JENIS_CRASH and tirukan_crash:
    return "terkirim, worker mati sebelum ack"
  antrean.selesai(pesan["id"])
  return "terkirim dan di-ack"
\end{lstlisting}

Pekerjaan \texttt{kerja-0003} menirukan kejadian yang paling sering
melahirkan duplikat: \textit{webhook}-nya terkirim, pekerjaannya ditandai
selesai, lalu \textit{worker}-nya mati sebelum meng-\textit{ack}. Tabel
\ref{tab:atleastonce} membandingkan jumlah percobaan dengan jumlah pengiriman
yang benar-benar tiba.

\begin{table}[!htb]
  \centering
  \caption{Percobaan \textit{worker} lawan pengiriman yang tiba di
           \textit{webhook}, hasil pengukuran nyata}
  \label{tab:atleastonce}
  \begin{tabularx}{\textwidth}{Xlll}
    \hline
    \textbf{Pekerjaan} & \textbf{Percobaan} & \textbf{Tiba} & \textbf{Akibat di tujuan} \\
    \hline
    \texttt{kerja-0001}, lancar & 1 & 1 & satu kali \\
    \hline
    \texttt{kerja-0002}, ditolak dua kali & 3 & 3 & satu kali, pada percobaan ketiga \\
    \hline
    \texttt{kerja-0003}, \textit{worker} mati & 2 & 1 & satu kali \\
    \hline
    \textbf{Jumlah} & \textbf{6} & \textbf{5} & \textbf{3 pekerjaan tuntas} \\
    \hline
  \end{tabularx}
\end{table}

Baris ketiga memuat hasil yang dicari. Pekerjaannya dikirim \textit{worker}
dua kali, tetapi hanya satu yang sampai ke \textit{webhook}, karena
pengiriman kedua berhenti di pemeriksaan duplikat. Tanpa pemeriksaan itu,
tujuan menerima pekerjaan yang sama dua kali, dan pada pekerjaan seperti
pembayaran akibatnya bukan hal kecil.

Kuncinya adalah id pekerjaan yang tetap, bukan id yang dibuat ulang tiap
percobaan. Dua percobaan \texttt{kerja-0003} membawa id yang sama, sehingga
dapat dikenali sebagai pekerjaan yang sama. Bila idnya dibuat baru tiap
pengiriman, keduanya terlihat seperti dua pekerjaan berbeda, dan tidak ada
yang dapat menahannya.

\section{Scheduled Job dan Eventual Consistency}

\textit{Scheduled job} adalah pekerjaan yang dijalankan menurut waktu, bukan
menurut permintaan. Pada demo ini, putaran \textit{worker} yang berulang tiap
0,25 detik berperan seperti itu: ia memeriksa antrean, mengambil pekerjaan
yang sudah waktunya, lalu tidur lagi. Pola yang sama dipakai pekerjaan harian
seperti merangkum laporan atau membersihkan data lama.

Akibat dari alur asinkron ini perlu dinyatakan terus terang. Begitu pekerjaan
dititipkan ke antrean, keadaan sistem tidak lagi seragam di satu waktu. Tabel
\ref{tab:konsisten} memuat urutannya dari hasil pengukuran, dan di sana
terlihat satu sistem sudah tahu sedangkan yang lain belum.

\begin{table}[!htb]
  \centering
  \caption{Keadaan kedua sistem pada beberapa titik waktu, hasil pengukuran
           nyata}
  \label{tab:konsisten}
  \begin{tabularx}{\textwidth}{lXX}
    \hline
    \textbf{Detik} & \textbf{Sistem kita} & \textbf{Tujuan \textit{webhook}} \\
    \hline
    0,00 & tiga pekerjaan diterima & belum menerima apa pun \\
    \hline
    0,01 & dua tuntas, satu gagal & \texttt{kerja-0001} dan \texttt{kerja-0003} \\
    \hline
    0,51 & duplikat ditahan & masih belum tahu \texttt{kerja-0002} \\
    \hline
    1,52 & tiga tuntas & ketiganya sudah diterima \\
    \hline
  \end{tabularx}
\end{table}

Keadaan seperti itu disebut \textit{eventual consistency}: kedua sisi akan
sama, tetapi tidak sekarang. Dua akibat praktisnya perlu disiapkan lebih
dahulu. Pertama, tampilan tidak boleh berjanji lebih daripada yang sudah
terjadi, sehingga kalimat "sedang diproses" lebih jujur daripada "sudah
terkirim". Kedua, jarak waktunya harus dipantau, karena antrean yang
menumpuk adalah masalah yang tidak terlihat pada halaman mana pun.

\section{Ringkasan}

Tiga transport membawa peristiwa yang sama dengan biaya yang berbeda, dan
selisihnya terukur. \textit{Polling} dengan jeda 0,5 detik menyerahkan
peristiwa dalam 199,8 milidetik dan memakai sebelas permintaan, sedangkan SSE
dan WebSocket menyerahkannya dalam 1,8 dan 1,9 milidetik lewat satu
sambungan. Keterlambatan \textit{polling} ditentukan jedanya, bukan kecepatan
servernya.

Pilihan di antara SSE dan WebSocket tidak ditentukan angka di tabel itu. SSE
berjalan satu arah, berbentuk teks, dan menyambung ulang sendiri di browser.
WebSocket berjalan dua arah dan lebih hemat pada pesan kecil, tetapi
penyambungan ulangnya ditulis sendiri. Papan pemberitahuan cukup dengan yang
pertama, sedangkan penyuntingan bersama membutuhkan yang kedua.

Alur asinkron memindahkan pekerjaan lama keluar dari permintaan, dan
memindahkan masalahnya juga. Antrean pada demo ini menggeser waktu siap tiap
pekerjaan yang diambil, sehingga pekerjaan tidak hilang ketika
\textit{worker}-nya mati, dengan ganti bahwa pekerjaan dapat terkirim lebih
dari sekali.

Pengukurannya memperlihatkan keduanya sekaligus. Enam percobaan
\textit{worker} menghasilkan lima pengiriman yang tiba, dan tiga pekerjaan
tuntas. Pekerjaan yang ditolak dua kali dicoba ulang pada detik 0,51 dan
1,52, yaitu jeda dasar lalu dua kali jeda dasar. Pekerjaan yang
\textit{worker}-nya mati sesudah mengirim dikirim ulang, lalu ditahan
pemeriksaan duplikat sebelum menyentuh tujuannya.

Pesan utama bab ini adalah tempat jaminan itu ditegakkan. Antrean hanya dapat
menjamin pekerjaan tidak hilang. Jaminan bahwa akibatnya tidak terjadi dua
kali hanya dapat ditegakkan \textit{consumer}, dengan id pekerjaan yang tetap
dan satu pemeriksaan yang dilakukan sebelum pekerjaannya dikerjakan.

\section{Latihan}

Seluruh file yang dipakai pada latihan ini tersedia di
\berkas{sources/chapter13/}, dihitung relatif terhadap akar repositori.
Latihan ini hanya membutuhkan Python 3, tanpa Docker dan tanpa pustaka
tambahan, termasuk WebSocket-nya.

Seluruh permintaan diarahkan ke layanan yang berjalan di komputer sendiri,
dan sifatnya mengamati, bukan membebani. Angka milidetik tiap mahasiswa akan
berbeda, sedangkan pola perbandingannya dan jumlah percobaannya akan sama.

Penyiapannya tiga perintah, dan urutannya berpengaruh, seperti pada Listing
\ref{lst:c13-siapkan}.

\begin{lstlisting}[language=bash, caption={Menyalakan ketiga layanan, tersedia di \berkas{sources/chapter13/}}, label={lst:c13-siapkan}]
cd sources/chapter13
python3 aplikasi.py &
python3 webhook.py &
python3 websocket.py &

# Keluarannya:
# Melayani di http://127.0.0.1:8100/sse dan /polling
# Webhook di http://127.0.0.1:8102/webhook, dua percobaan pertama ditolak
# WebSocket di ws://127.0.0.1:8101/
\end{lstlisting}

\subsection*{Latihan 1: Membandingkan Tiga Transport}

Latihan ini mengukur keterlambatan dan biaya ketiga transport pada peristiwa
yang sama, lalu memisahkan biaya yang terlihat dari biaya yang tersembunyi.

\begin{enumerate}
  \item Jalankan \berkas{ukur-transport.py}, seperti pada Listing
        \ref{lst:c13-lat1-jalan}, lalu catat ketiga barisnya.
  \item Catat jumlah permintaan tiap transport, lalu hitung berapa permintaan
        yang jawabannya kosong pada varian \textit{polling}.
  \item Buka \texttt{http://127.0.0.1:8100/sse} dengan \texttt{curl -sN},
        lalu terbitkan satu peristiwa dari terminal lain. Catat bentuk blok
        yang tiba.
  \item Isi Tabel \ref{tab:c13-lat1-hasil}. Tabel \ref{tab:c13-lat1-contoh}
        memperlihatkan contoh yang sudah terisi.
  \item Jawab tiga pertanyaan berikut. Pertama, mengapa keterlambatan
        \textit{polling} tidak membaik walaupun servernya cepat? Kedua,
        mengapa kolom byte tidak cukup untuk menilai biaya \textit{polling}?
        Ketiga, kapan WebSocket pantas dipilih walaupun SSE lebih sederhana?
\end{enumerate}

\begin{lstlisting}[language=bash, caption={Mengukur ketiga transport, tersedia di \berkas{sources/chapter13/ukur-transport.py}}, label={lst:c13-lat1-jalan}]
python3 ukur-transport.py

# Keluarannya:
#   Transport   Peristiwa  Tengah ms  Tertinggi     Byte  Permintaan
#   polling            10      199.8      202.0      639          11
#   sse                10        1.8        2.0      400           1
#   websocket          10        1.9        2.1      490           1
\end{lstlisting}

\begin{table}[!htb]
  \centering
  \caption{Contoh hasil Latihan 1 yang sudah terisi}
  \label{tab:c13-lat1-contoh}
  \begin{tabularx}{\textwidth}{Xlll}
    \hline
    \textbf{Transport} & \textbf{Tengah ms} & \textbf{Byte} & \textbf{Permintaan} \\
    \hline
    \textit{polling} 0,5 detik & 199,8 & 639 & 11 \\
    \hline
    SSE & 1,8 & 400 & 1 \\
    \hline
    WebSocket & 1,9 & 490 & 1 \\
    \hline
  \end{tabularx}
\end{table}

\begin{table}[!htb]
  \centering
  \caption{Lembar isian hasil pengukuran Latihan 1}
  \label{tab:c13-lat1-hasil}
  \begin{tabularx}{\textwidth}{Xlll}
    \hline
    \textbf{Transport} & \textbf{Tengah ms} & \textbf{Byte} & \textbf{Permintaan} \\
    \hline
    \textit{polling} 0,5 detik & \dots & \dots & \dots \\
    \hline
    SSE & \dots & \dots & \dots \\
    \hline
    WebSocket & \dots & \dots & \dots \\
    \hline
  \end{tabularx}
\end{table}

\subsection*{Latihan 2: Jeda Polling dan Keterlambatannya}

Latihan ini membuktikan bahwa keterlambatan \textit{polling} ditentukan
jedanya, dan bahwa memperkecil jeda membayar dengan jumlah permintaan.

\begin{enumerate}
  \item Ubah \texttt{JEDA\_POLLING\_DETIK} di \berkas{ukur-transport.py}
        menjadi 0,1 lalu jalankan pengukurannya. Catat keterlambatan dan
        jumlah permintaannya.
  \item Ulangi dengan nilai 0,5 dan 1,0.
  \item Bandingkan ketiganya dengan selisih jeda penerbitan peristiwa, yaitu
        \texttt{JEDA\_ANTARPENERBITAN\_DETIK} yang bernilai 0,3 detik.
  \item Isi Tabel \ref{tab:c13-lat2-hasil}. Tabel \ref{tab:c13-lat2-contoh}
        memperlihatkan contoh yang sudah terisi.
  \item Jawab dua pertanyaan berikut. Pertama, mengapa keterlambatannya
        mendekati jeda \textit{polling} dikurangi jeda penerbitan, bukan
        setengah jeda \textit{polling}? Kedua, pada beban nyata dengan
        peristiwa yang datang sembarang waktu, berapa keterlambatan rata-rata
        yang pantas diharapkan?
\end{enumerate}

\begin{lstlisting}[language=bash, caption={Mengukur satu transport pada beberapa jeda, memakai \berkas{sources/chapter13/ukur-transport.py}}, label={lst:c13-lat2-jalan}]
# JEDA_POLLING_DETIK diubah, lalu pengukurannya dijalankan ulang:
python3 ukur-transport.py

# Baris polling pada tiga nilai jeda, 0.1 lalu 0.5 lalu 1.0 detik:
#   polling            10        3.3        5.3      997          31
#   polling            10      200.2      202.6      657          11
#   polling            10      700.1      701.0      657          11
\end{lstlisting}

\begin{table}[!htb]
  \centering
  \caption{Contoh hasil Latihan 2 yang sudah terisi}
  \label{tab:c13-lat2-contoh}
  \begin{tabularx}{\textwidth}{Xlll}
    \hline
    \textbf{Jeda polling} & \textbf{Tengah ms} & \textbf{Byte} & \textbf{Permintaan} \\
    \hline
    0,1 detik & 3,3 & 997 & 31 \\
    \hline
    0,5 detik & 200,2 & 657 & 11 \\
    \hline
    1,0 detik & 700,1 & 657 & 11 \\
    \hline
  \end{tabularx}
\end{table}

\begin{table}[!htb]
  \centering
  \caption{Lembar isian hasil pengukuran Latihan 2}
  \label{tab:c13-lat2-hasil}
  \begin{tabularx}{\textwidth}{Xlll}
    \hline
    \textbf{Jeda polling} & \textbf{Tengah ms} & \textbf{Byte} & \textbf{Permintaan} \\
    \hline
    0,1 detik & \dots & \dots & \dots \\
    \hline
    0,5 detik & \dots & \dots & \dots \\
    \hline
    1,0 detik & \dots & \dots & \dots \\
    \hline
  \end{tabularx}
\end{table}

\subsection*{Latihan 3: Retry dan Jeda yang Berlipat}

Latihan ini memperlihatkan jeda \textit{retry} bekerja, dan membuktikan
angkanya dapat dihitung sebelum diukur.

\begin{enumerate}
  \item Jalankan \berkas{ukur-antrean.py}, seperti pada Listing
        \ref{lst:c13-lat3-jalan}, lalu catat keenam percobaannya beserta
        detiknya.
  \item Hitung selisih detik antarpercobaan \texttt{kerja-0002}, lalu
        bandingkan dengan \texttt{JEDA\_DASAR\_DETIK} di
        \berkas{antrean.py}.
  \item Ubah \texttt{JEDA\_DASAR\_DETIK} menjadi 1,0 lalu ulangi, dan catat
        detik percobaan kedua dan ketiganya.
  \item Isi Tabel \ref{tab:c13-lat3-hasil}. Tabel \ref{tab:c13-lat3-contoh}
        memperlihatkan contoh yang sudah terisi.
  \item Jawab tiga pertanyaan berikut. Pertama, berapa lama satu pekerjaan
        bertahan sebelum percobaan kelimanya, bila jeda dasarnya 0,5 detik?
        Kedua, mengapa \textit{retry} tanpa jeda memperburuk keadaan? Ketiga,
        ke mana pekerjaan yang tetap gagal sesudah batas percobaan pantas
        dipindahkan?
\end{enumerate}

\begin{lstlisting}[language=bash, caption={Menjalankan satu daur antrean, tersedia di \berkas{sources/chapter13/ukur-antrean.py}}, label={lst:c13-lat3-jalan}]
python3 ukur-antrean.py

# Keluarannya:
#   Pekerjaan    Jenis                 Ke   Detik  Hasil
#   kerja-0001   kirim-laporan          1    0.01  terkirim dan di-ack
#   kerja-0002   gagal-dua-kali         1    0.01  gagal, webhook 500
#   kerja-0003   crash-setelah-kirim    1    0.01  terkirim, mati sebelum ack
#   kerja-0002   gagal-dua-kali         2    0.51  gagal, webhook 500
#   kerja-0003   crash-setelah-kirim    2    0.51  duplikat ditahan
#   kerja-0002   gagal-dua-kali         3    1.52  terkirim dan di-ack
\end{lstlisting}

\begin{table}[!htb]
  \centering
  \caption{Contoh hasil Latihan 3 yang sudah terisi}
  \label{tab:c13-lat3-contoh}
  \begin{tabularx}{\textwidth}{Xlll}
    \hline
    \textbf{Jeda dasar} & \textbf{Percobaan 2} & \textbf{Percobaan 3} & \textbf{Selisih} \\
    \hline
    0,5 detik & 0,51 & 1,52 & 0,5 lalu 1,0 \\
    \hline
  \end{tabularx}
\end{table}

\begin{table}[!htb]
  \centering
  \caption{Lembar isian hasil pengukuran Latihan 3}
  \label{tab:c13-lat3-hasil}
  \begin{tabularx}{\textwidth}{Xlll}
    \hline
    \textbf{Jeda dasar} & \textbf{Percobaan 2} & \textbf{Percobaan 3} & \textbf{Selisih} \\
    \hline
    0,5 detik & \dots & \dots & \dots \\
    \hline
    1,0 detik & \dots & \dots & \dots \\
    \hline
  \end{tabularx}
\end{table}

\subsection*{Latihan 4: Duplikat dan Consumer Idempoten}

Latihan ini membuktikan bahwa pengiriman dapat terjadi dua kali, dan bahwa
akibat gandanya hanya dapat ditahan \textit{consumer}.

\begin{enumerate}
  \item Dari hasil Latihan 3, catat jumlah percobaan \textit{worker} dan
        jumlah pengiriman yang tiba di \textit{webhook}.
  \item Catat pekerjaan mana yang dikirim dua kali oleh \textit{worker}
        tetapi hanya sekali tiba di \textit{webhook}.
  \item Matikan pemeriksaan duplikat dengan memberi komentar pada pemanggilan
        \texttt{sudah\_diproses} di \berkas{worker.py}, jalankan ulang, lalu
        catat jumlah pengiriman yang tiba sekarang.
  \item Pulihkan kembali pemeriksaannya.
  \item Isi Tabel \ref{tab:c13-lat4-hasil}. Tabel \ref{tab:c13-lat4-contoh}
        memperlihatkan contoh yang sudah terisi.
  \item Jawab tiga pertanyaan berikut. Pertama, mengapa id pekerjaan harus
        tetap pada tiap percobaan? Kedua, mengapa pemeriksaan duplikat
        diletakkan sebelum pengiriman, bukan sesudahnya? Ketiga, pekerjaan
        apa saja yang akibatnya memang tidak berubah walaupun dikerjakan dua
        kali?
\end{enumerate}

\begin{lstlisting}[language=bash, caption={Membandingkan percobaan dengan pengiriman yang tiba, tersedia di \berkas{sources/chapter13/ukur-antrean.py}}, label={lst:c13-lat4-jalan}]
python3 ukur-antrean.py | tail -3

# Keluarannya:
#   Percobaan worker: 6, pengiriman tiba: 5, pekerjaan tersisa: 0,
#   ditandai diproses: 3
\end{lstlisting}

\begin{table}[!htb]
  \centering
  \caption{Contoh hasil Latihan 4 yang sudah terisi}
  \label{tab:c13-lat4-contoh}
  \begin{tabularx}{\textwidth}{Xlll}
    \hline
    \textbf{Keadaan} & \textbf{Percobaan} & \textbf{Tiba} & \textbf{Akibat ganda} \\
    \hline
    pemeriksaan duplikat hidup & 6 & 5 & tidak ada \\
    \hline
    \texttt{kerja-0003} saja & 2 & 1 & tidak ada \\
    \hline
  \end{tabularx}
\end{table}

\begin{table}[!htb]
  \centering
  \caption{Lembar isian hasil pengamatan Latihan 4}
  \label{tab:c13-lat4-hasil}
  \begin{tabularx}{\textwidth}{Xlll}
    \hline
    \textbf{Keadaan} & \textbf{Percobaan} & \textbf{Tiba} & \textbf{Akibat ganda} \\
    \hline
    pemeriksaan duplikat hidup & \dots & \dots & \dots \\
    \hline
    pemeriksaan duplikat mati & \dots & \dots & \dots \\
    \hline
  \end{tabularx}
\end{table}
